\chapter{Анализ предметной области}
\section{Классификация изображений растений}
Королевские ботанические сады в Кью опубликовали Всемирный список сосудистых растений, в котором задокументировано 350 386 признанных видов сосудистых растений~\cite{ref1102}. Автоматизация классификации такого огромного множества видов~---~сложная задача, которая стала возможной благодаря машинному обучению и методам компьютерного зрения.

Классическая задача идентификации растений основана на ручной работе специалистов-таксономистов и требует значительных временных и экспертных ресурсов. С увеличением количества доступных визуальных данных растёт потребность в автоматизированной классификации растений~\cite{ref1101}. 

\section{Особенности визуальных данных}
 Незначительные различия в форме, текстуре или цвете листьев могут разделять близкородственные виды. Сложность классификации усугубляется внешними факторами, такими как различия в фоне, позе и освещении, которые приводят к значительным различиям внутри классов изображений~\cite{ref1103}. Эта проблема решается за счёт извлечения дискриминационных признаков, которые минимизируют внутриклассовые различия и максимально увеличивают межклассовые различия.
 
Традиционно для классификации растений анализируют изображения листьев, выделяя множество признаков, таких как текстура листьев, их форма и цвет. Цветовая характеристика в большей степени зависит от классификации и идентификации видов растений, поскольку цвет листьев может меняться в зависимости от окружающей среды в разные сезоны. Характеристики текстуры в большей степени основаны на информации, полученной из его жилок и жилковании. В последнее время узоры жилок листьев считаются важным фактором для идентификации видов растений, при этом существует лишь несколько методов извлечения структуры жилок листьев~\cite{ref1101}. 

Помимо проблемы разнообразия видов и подвидов растений существует проблема количества имеющихся данных. Например, на наборе PlantCLEF 2019, сосредоточенном на тропических регионах с недостаточным количеством данных, лучшее решение имеет всего 41\% точности~\cite{ref1104} по сравнению с
96\% точностью в наборах данных по флоре умеренного пояса~\cite{ref1105}
